\section*{September 22, 2026: Removing foreclosure-related sales from the price sample}
\textit{Update, later September 22: REO resales are now flagged rather than removed; see ``Flagging REO resales and price tails instead of removing them.''}

\textit{Drafted by Claude; Jacob decided to drop these sales and chose the deed-type fallback.}

The clean price sample previously retained foreclosure auctions and resales by
lenders and government agencies. These are not market sales on comparable
terms, and published hedonic studies usually exclude them. In 2008--2013 they
were a large share of Chicago transactions: 29--45\% of sales that pass the
county flags each year, falling to 14\% by 2018. They sold for a median of
\$63,000 nominal, against \$230,000 for other sales. The citywide 2008--2018
sample falls from 167,468 to 124,396 sales. The median real price rises most
in the crash years. In 2009 it moves from \$181,000 to \$283,000, so the fall
from 2008 shrinks from 32\% to 18\%.

The rule reads buyer and seller names, which are reported in every year, so it
treats the pre- and post-2013 periods alike. It removes three kinds of sale.
The first is a foreclosure auction: a judicial-sale company, sheriff or selling
officer is the seller. The second is a resale by a lender, servicer,
securitization trustee, Fannie Mae, Freddie Mac, HUD, the VA or the Cook County
Land Bank. The third is a transfer to one of these. Banks acting as Chicago
land trustees are kept. Their sales have a median price near \$180,000 and
are ordinary transactions. Detailed MyDec deed types begin only in 2013, so
they cannot define the rule. Among post-2012 sales with a known seller,
though, they validate it. The names flag 98\% of judicial-sale deeds and 80\%
of special warranty deeds, the usual REO instrument. They flag under 1\% of
ordinary warranty deeds and 3\% of trustee deeds, the usual land-trust
instrument.

Seller names are missing for 46\% of 2015 sales. Using names alone, the rule
removed only 2.7\% of 2015 sales. When the seller name is missing, the rule
now falls back on the MyDec deed type. Judicial and court-officer deeds, deeds
in lieu of foreclosure and special warranty deeds count as
foreclosure-related. The 2015 removal share becomes 20.9\%, between 2014
(24.6\%) and 2016 (19.4\%). Names are almost never missing before 2014, so the
fallback does not change how the rule treats the pre-period.

In the half-mile single-family packet, the rule keeps 63\% of treated and 60\%
of control sales in 2008--2012. It keeps 78\% and 74\% in 2014--2018. Removal
is thus similar across groups but heavier before 2013. Removing these cheap
sales also raises the within-year 99.9th-percentile price-per-square-foot
cutoffs, so 26 previously trimmed sales at \$850--960 per square foot remain. The rule misses
foreclosure-related sales whose parties are named individuals, such as
selling officers, and unrecognized servicers. It also misses short sales,
which the public export does not identify. Later investor resales of
foreclosed homes remain; they are the next cleaning question.

Later the same day, a review of the lowest-priced remaining sales found REO
sellers with consumer-finance names, such as Household Finance, Accredited
Home Lenders, Residential Funding and Liquidation Properties. Adding these
names removes 948 more sales, with a median price of \$50,000. Broader tokens
such as FINANCE or FUNDING were rejected because they matched surnames,
relocation companies, and investor-financing firms. After this addition and
the unit-count change recorded separately, the citywide 2008--2018 sample is
126,498 sales.

Jacob also retired the September 7 descriptive sample that kept sales with
incomplete characteristics. The data-overview packets had rebuilt the cleaning
rules independently, so every cleaning change had to be made twice. They now
use the production clean file directly. \texttt{clean\_home\_sales} also writes a
2008--2023 file for the longer packet; its 2008--2018 rows equal the 2018 file.
The half-mile single-family packet falls from 11,206 to 7,611 sales (1,211
treated, 6,400 control); the quarter-mile all-property packet falls from
11,600 to 7,190. At a half mile, only 4 of the 11,206 earlier sales lacked
complete characteristics, so almost all of that fall is the foreclosure rule.
Regression audits have not been rerun.

\small
Reproduction: run \texttt{make} in \texttt{tasks/clean\_home\_sales/code}.
The patterns are literal in \texttt{clean\_home\_sales.R}; the build prints
annual removal counts. Existing parcel-sales snapshot; working tree based on
\texttt{e76f4dd}.
\normalsize
